\documentclass[../main.tex]{subfiles}
\begin{document}

\section{Pasarela}
La pasarela (Gateway) que se utilizó en el proyecto, se construyo a partir del concentrador \gls{lwan} del fabricante IMST, del modelo  referencia  "WiMOD iC880A". El host utilizado, para  conversión de protocolos y la comunicación con la nube, fue una plaqueta de computación  Raspberry pi-3.
\begin{figure}[h]
    \centering
    \includegraphics{gw/DIagrama_Lora.png}
    \caption{Diagrama de la pasarela del proyecto}
    \label{fig:bloque}
\end{figure}

\subsection{Concentrador}
El concentrador se comunica con los nodos asociados a este, y hace la recepción de los datos en los formatops y protocolos \gls{lwan}, las especificaciones del concentrador: \\

\begin{tabular}{p{7cm}l}
    \hline
    {\bf Especificaciones} & {\bf Valor} \\
    \hline
    Banda de frecuencias  &   868 MHz \\
    Sensibilidad de bajada & -137 dBm\\
    Interfase & SPI\\
    Alimentación & 5 V\\
    Potencia de subida & 20 dBm\\
    Rango en linea de vista & 15 km\\
    Rango en zona urbana & algunos km \\ \hline
\end{tabular} \\

El concentrador  "WiMOD iC880A", incluye el procesador banda base "Semtech SX1301" , el cual es un procesador digital para la banda \gls{lwan}, el cual recibe el microcodiog del host, cuando la comunicación retorna hacia el nodo. Contien 10 canales de recepción, lo cual permite varias formas de programación en la recepción de los nodos.  
El modulo de radio, tiene los dos modos de recpción: la modulación LoRa y la modulación GFSK o FSK,  incluye el modulo de radio. 

\begin{figure}
    \caption{Plaqueta del concentrador "Semtech SX1301"}
    \centering
    \includegraphics[width=0.8\textwidth]{gw/concentrador.png}
    \label{fig:my_label}
\end{figure}

\end{document}